\chapter{Modelo matemático}
\section{Planta térmica}
Desde el punto de vista termodinámico es conveniente pensar al APD como una red de nodos térmicos que se conectan mediante resistencias térmicas.\\
Los componentes termicos de la planta del sistema esta constituidos principalmente por una celda Peltier, un termistor, un disipador de energía y el ambiente, el objetivo principal es bajar y estabilizar la temperatura del sistema en un rango efectivo de 0° a 10°C mediante la cara fría de la TEC y la dicipación del calor de la cara caliente hacia el ambiente. Esto aumenta la capacidad de respuesta de la región activa del APD.
(colocar una imagen de la red de nodos)\\
Se comienza por modelar la energia interna del sistema mediante la primera ley de la termodinámica 
\begin{equation}
    dU = {d}Q + {d}W
\end{equation}
el trabajo mecánico en el sistema es nulo, la energia interna del sistema esta solo en terminos del calor neto transferido 
\begin{equation}
     dU = {d}Q_{neto}
\end{equation}
en donde 
\begin{equation}
    Q_{neto} = Q_{in}-Q_{out}
\end{equation}
entonces la ecuación (2.3) puede escribirse como
\begin{equation}
    dU = dQ_{in}-dQ_{out}
\end{equation}
En el modelo $Q_{in}$ corresponde al flujo de calor suministrado por la TEC al sistema $Q_{TEC}$ y $Q_{out}$ corresponde al calor que disipa el sistema al ambiente $Q_{APD}$ por lo que se puede escribir 
\begin{equation}
    dU = dQ_{TEC}-dQ_{APD}
\end{equation}
se considero al tiempo t y a la temperatura T como las principales variables termodinámicas asociadas a la energia interna del sistema 
\begin{equation}
    U = U(T,t)
\end{equation}
se toma la derivada total en términos de un parámetro s 
\begin{equation}
    dU(T,t) = (\frac{\partial U}{\partial T})_{t} \frac{dT}{ds} + (\frac{\partial U}{\partial t})_{T}\frac{dt}{ds}
\end{equation}
como se estabilizara la temperatura podemos anular las variaciones de la misma 
\begin{equation}
    dU = (\frac{\partial U}{\partial t})_{T}\frac{dt}{ds}
\end{equation}
Para el calor se considero a la distancia x y el tiempo t como las principales variables termodinámicas 
\begin{equation}
    Q = Q(x,t)
\end{equation}
se relaciona al calor con la temperatura mediante la capacidad térmica del sistema C
\begin{equation}
    Q(x,t)=CdT(x,t)
\end{equation}
se toma la derivada total respecto al parametro s 
\begin{equation}
    dT(x,t) = C(\frac{\partial T}{\partial x})_{t}\frac{dx}{ds} + C(\frac{\partial T}{\partial t})_{x}\frac{dt}{ds}
\end{equation}
como las dimensiones del sistema no cambian, la variacion de la distancia es nula 
\begin{equation}
    dQ = C(\frac{\partial T}{\partial t})_{x}\frac{dt}{ds}
\end{equation}
sustituyendo (2.12) y (2.8) en (2.2) se consigue
\begin{equation}
     (\frac{\partial U}{\partial t})_{T}\frac{dt}{ds} =
     C(\frac{\partial T}{\partial t})_{x}\frac{dt}{ds}
\end{equation}
esta igualdad puede escribirse en terminos de derivadas totales 
\begin{equation}
    \frac{dU}{dt} = C\frac{dT}{dt}
\end{equation}
combinando este resultado con (2.5) se consigue
\begin{equation}
    C\frac{dT(t)}{dt} = \frac{dQ_{TEC}}{dt}-\frac{dQ_{APD}}{dt}
\end{equation}
\begin{equation}
    C\frac{dT(t)}{dt} = \dot{Q}_{TEC}-\dot{Q}_{APD}
\end{equation}
\subsection{Calor absorbido por la TEC}
El calor absorbido por la cara fría de la TEC se modela mediante una función que depende de la temperatura de la cara caliente $T_{h}$, la cara temperatura de la cara fría $T_{c}$ y la corriente suministrada $I$
\begin{equation}
    \dot{Q}_{TEC}(T_{h}, T_{c}, I)= \alpha T_{c}I-\frac{1}{2}I^{2}R-K_{TEC}(T_h-T_c)
\end{equation}
se considera un punto $(T_0c, T_0h, I_0)$ cerca del (0,0,0) y se toma la serie de Taylor al rededor de el, cortando los términos a primer orden se tiene 
\begin{equation}
    \begin{split}
    \dot{Q}_{TEC}(T_{c}(t), T_{h}(t),I(t))= \alpha T_{0c}I_{0} -\frac{1}{2}I_{0}^2R -K_{TEC}(T_{0h}-T_{0c}) +(\alpha I_{0}+K_{TEC})(T_{c}(t)-T_{0c})\\ 
    -K_{TEC}(T_{h}(t)-T_{oh}) - (\alpha T_{0c}-I_{0}R)(I(t)-I_{0})
    \end{split}
\end{equation}
esta ecuación se simplifica al considerar que $(T_{0c}, T_{oh},I_{0})$ esta cerca del cero. Se multiplica por $-1$ para indicar que el sistema esta enfriando 
\begin{equation}
     \dot{Q}_{TEC}(T_{c}(t), T_{h}(t),I(t))= -(\alpha I_{0} + K_{TEC} )\Delta T_{c}(t)+K_{TEC}\Delta T_{h}(t) - (\alpha T_{0c}-I_{0}R)\Delta I(t)
\end{equation}

la corriente se suministra solo al momento de encender la planta y la acción no es inmediata, así que lo más adecuado es modelar la inyección de corriente mediante una función $u(t)$ de Heavisade con un retardo L en el tiempo
\begin{equation}
\Delta  I(t-L)u(t-L) =
\begin{cases}
    \Delta I(t-L) & \text{si } t \geq L\\
    0  & \text{si } t < L
\end{cases}
\end{equation}
haciendo $\gamma =(\alpha I_{0} + K_{TEC} )$ y $b=(\alpha T_{0c}-I_{0}R)$ se consigue 
\begin{equation}
     \dot{Q}_{TEC}(T_{c}(t), T_{h}(t),I(t))= -\gamma \Delta T_{c}(t)+K_{TEC}\Delta T_{h}(t) - b\Delta I(t-L)u(t-L)
\end{equation}
\subsection{Intercambio de calor con el ambiente}
Para modelar el intercambio de calor entre el sistema y el ambiente se emplea la ecuación de difusión de 
\begin{equation}
    \dot{Q}_{APD} = -k\nabla T
\end{equation}
al estar establecida una dirección para el  flujo de calor se puede escribir 
\begin{equation}
    \dot{Q}_{APD}(x,t)=-k\frac{\partial T(x,t)}{\partial x}
\end{equation}
Sea A el área del disipador de energía térmica, se define entonces el calor total emitido 
\begin{equation}
     \dot{Q}_{APD}(x,t)=-kA\frac{\partial T(x,t)}{\partial x}
\end{equation}
sea r la distancia entre la cara caliente de la TEC y el extremo exterior del disipador de energía, se puede escribir entonces 
\begin{equation}
    \dot{Q}_{APD}(x,t)=kA\frac{T(r,t)-T(0,t)}{r}
\end{equation}
haciendo que $T(r,t)=T_{amb}(t)$, $T(0,t)=T_{APD}(t)$ y $R_{term}=\frac{r}{kA}$ se puede escribir 
\begin{equation}
     \dot{Q}_{APD}(x,t)=-\frac{1}{R_{term}}(T_{amb}(t)-T_{APD}(t))
\end{equation}
sustituyendo (2.26) y (2.21) en (2.16) se consigue
\begin{equation}
    C\frac{dT(t)}{dt} =  - \gamma \Delta T_{c}(t)+K_{TEC}\Delta T_{h}(t) - b\Delta I(t-L)u(t-L) + \frac{1}{R_{term}}(T_{amb}(t)-T_{APD}(t))
\end{equation}
se busca poner todo en terminos de cambios y que  
\begin{equation}
    T(t)=T_{c}(t)=T_{APD}(t)
\end{equation}
para la retroalimentacion negativa se define 
\begin{equation}
    \Delta T_(t) = T(t)-T_{oc}
\end{equation}
tomando la derivada respecto al tiempo se tiene 
\begin{equation}
    \frac{d\Delta T(t)}{dt} = \frac{dT(t)}{dt}
\end{equation}
por otro lado, debido a que $T_{oc}= 0$ se puede escribir 
\begin{equation}
    \Delta T_{APD}(t) = T_{APD}(t)-T_{oc} = T_{APD}(t)
\end{equation}
aplicando esta misma idea a la temperatura ambiente se tiene 
\begin{equation}
    \Delta T_{amb}(t) = T_{amb}(t)-T_{oc} = T_{amb}(t)
\end{equation}
al aplicar esto y agrupar se consigue la ecuación diferencial de primer orden que modela a la planta térmica del APD como un sistema termodinámico de enfriamiento. 
\begin{equation}
    C\frac{d\Delta T(t)}{dt} = -a\Delta T(t) + K_{TEC}\Delta T_{h}(t) - b\Delta I(t-L)u(t-L)+ \frac{1}{R_{term}}\Delta T_{amb}(t)
\end{equation}
siendo $a = \gamma+\frac{1}{R_{term}}$
\section{Funciones de transferencia}
Para cualquier sistema cuya relación de entrada $E(t)$ y salida S(t) pueda modelarese mediante  ecuaciones diferenciales lineales se define su función de transferencia como 
\begin{equation}
    G(s) = \frac{\Lap{S(t)}}{\Lap{E(t)}} = \frac{S(s)}{E(s)}
\end{equation}
Se considera la definición de la transformada de Laplace:
La transformada de Laplace es una herramienta matemática que toma una función en el tiempo y la convierte en otra función de variable compleja s.
\begin{equation}
    F(s)=\Lap {f(t)}=\int_{0}^{\infty}f(t)e^{-st}dt
\end{equation}
se aplica a la ecuación diferencial 
\begin{equation}
    \Lap{C\frac{d\Delta T(t)}{dt}} =\Lap{-a\Delta T (t) + K_{TEC}\Delta T_{h}(t) - b\Delta I(t-L)u(t-L)+ \frac{1 }{R_{term}}\Delta T_{amb}(t)}
\end{equation}
\begin{equation}
    C\Lap{\frac{d\Delta T(t)}{dt}} = -a\Lap{\ \Delta T (t)} + K_{TEC}\Lap{\Delta T_{h(t)} - b\Lap{\Delta I(t-L)u(t-L)}+ \frac{1 }{R_{term}}\Lap{\Delta T_{amb}(t)}}
\end{equation}
\begin{equation}
    C(s\Delta T(s)-\Delta T(0)) = -a\Delta T(s) + K_{TEC}\Delta T_{h}(s)-be^{-Ls}\Delta I(s) + \frac{1}{R_{term}}\Delta T_{amb}(s)
\end{equation}
considerando $\Delta T(0)=\Delta T_{amb}(s)$ y agrupando se consigue 
\begin{equation}
    (Cs+a)\Delta T(s) = K_{TEC}\Delta T_{h}(s)-be^{-Ls}\Delta I(s) + (C+\frac{1}{R_{term}})\Delta T_{amb}(s)
\end{equation}
naturalmente la temperatura del sistema no depende solo de la corriente suministrada, depende de la temperatura de la cara caliente de la TEC y la temperatura del ambiente. Se puden establecer tres funcione de transferencia a primer orden;
\begin{itemize}
    \item una para  la corriente suministrada $T_{h}(s) = T_{amb}(s) = 0$  
    \begin{equation}
        \frac{\Delta T(s)}{\Delta I(s)} = \frac{-be^{-Ls}}{Cs+a}
    \end{equation}
    \item una para la cara caliente de la TEC $I(s) = T_{amb} = 0$  
    \begin{equation}
        \frac{\Delta T(s)}{\Delta T_{h}(s)} = \frac{K_{TEC}}{Cs+a}
    \end{equation}
    \item y una para el ambiente $T_{h}(s) = I(s)=0$
    \begin{equation}
        \frac{\Delta T(s)}{\Delta T_{amb}(s)} = \frac{CR_{term}+1}{R_{term}(Cs+a)}
    \end{equation}
    \end{itemize}
se emplea (2.32) para caracterizar  el enfriamiento del sistema a partir de la inyección de corriente, factorizando a se consigue 
\begin{equation}
     \frac{\Delta T(s)}{\Delta I(s)} = \frac{\frac{-b}{a}e^{-Ls}}{\frac{C}{a} s + 1}
\end{equation}
haciendo $K=\frac{b}{a} $ y $\tau=\frac{C}{a}$ se obtiene la forma  más común asociada a dicha función de transferencia de calor por suministro  de corriente  
\begin{equation}
      G(s)= \frac{-Ke^{-Ls}}{\tau s +1}
\end{equation}
se demuestra que $\tau$ tiene unidades de tiempo
\begin{equation}
    [\tau] = [\frac{C}{a}] = [\frac{C}{\gamma+\frac{1}{R_{term}}} ]
\end{equation}
la capacidad termica C tiene unidades de $\frac{J}{K}$, la resistencia térmica de $R_{term}$ tiene unidades de $\frac{K}{W}$ y $\gamma$ por su parte 
\begin{equation}
    [\gamma] = [\alpha I_{0} + K_{TEC}] = \frac{W}{K}
\end{equation}
se tiene  
\begin{equation}
    [\tau] = \frac{\frac{J}{K}}{\frac{W}{k}}=\frac{J}{W} = s
\end{equation}
\subsection{Conexión con el experimento}
Para conectar el modelo matemático con el experimento se debe pensar en el efecto de un escalón de corriente 
\begin{equation}
    \Delta I(t-L)u(t-L)=\Delta I_{0}u(t-L)
\end{equation} 
cuya transformada de Lapace es 
\begin{equation}
    \Delta I(s)= \frac{\Delta I_{0}e^{-sL}}{s}
\end{equation}
sustituyendo en (2.36) y despejando T(s) se consigue  
\begin{equation}
    \Delta T(s) = \frac{\Delta I_{0}Ke^{-sL}}{s(\tau s +1)}
\end{equation}
se descomponen los dos factores del denominador mediante el metodo de fracciones parciales
\begin{equation}
    \Delta T(s) = e^{-sL}(\frac{\Delta I_{0}K}{s} - \frac{\tau \Delta I_{0}K}{\tau s +1})
\end{equation}
se aplica ahora la transformada inversa de Laplce para analizar la respuesta en el tiempo del sistema modelado 
\begin{equation}
   \Lap{ \Delta T(s) }= \Lap {e^{-sL}(\frac{\Delta I_{0}K}{s} - \frac{\tau \Delta I_{0}K}{\tau s +1})}
\end{equation}
\begin{equation}
     \Lap{ \Delta T(s) } = \Delta I_{0}K \Lap{\frac{e^{-sL}}{S}} - \Delta I_{0}K \Lap{\frac{1}{s+\frac{1}{\tau}}}
\end{equation}
\begin{equation}
    \Delta T(t) = \Delta I_{0}Ku(t-L)(1-e^{-\frac{t}{\tau}})
\end{equation}
 Para $t \geq L$ se tiene 
 \begin{equation}
    \Delta T(t) = \Delta I_{0}K(1-e^{-\frac{t}{\tau}})
\end{equation}
cuand $t \xrightarrow{}\infty$ el sistema llega a la temperatura de equilibrio 
\begin{equation}
    \Delta T_{ss} =\Delta I_{0}K
\end{equation}
la ecuación (2.44) nos permite relacionar las variaciones de la temperatura del sistema con la variaciones de la   escalon de corriente suministrado.
\begin{equation}
    \Delta T(t) =\Delta I_{0}K(1-e^{-\frac{t}{\tau}})
\end{equation}
cuando $t \xrightarrow{}\infty$ el sistema tiende al equilibrio, en ese entonces 
\begin{equation}
     \Delta T_{ss}(t) =\Delta I_{0}K
\end{equation}
se define entonces la ganancia termica del sistema K en términos del cambio total de temperatura del sistema respecto al cambio total de corriente suministrada 
\begin{equation}
    K = \frac{\Delta T(t)}{\Delta I_{0}}=\frac{\Delta T_{ss}}{\Delta I_{0}}
\end{equation}
considerando ahora el cociente entre (2.44) y (2.45) se tiene 
\begin{equation}
    \frac{\Delta T(t)}{\Delta T_{ss}}=1-e^{-\frac{t}{\tau}}
\end{equation}
haciendo $t=\tau$ y considerando una pequeñua variación  
\begin{equation}
     \frac{\Delta T(\tau)}{\Delta T_{ss}}=1-e^{-1}=0.632121
\end{equation}

un cambio en la temperatura puede escribirse como 
\begin{equation}
    \Delta T(t)= T(t)-T_{amb}
\end{equation}
de modo tal que 
\begin{equation}
    T(\tau)= \Delta T(\tau)+ T_{amb}(\tau)
\end{equation}
se puede escribir entonces 
\begin{equation}
    T(\tau)= T_{amb}(\tau) + 0.632\Delta T_{ss}(\tau) 
\end{equation}
considerando que 
\begin{equation}
    \Delta T_{ss}(\tau)= T_{ss}(\tau) - T_{amb}(\tau)
\end{equation}
se obtiene 
\begin{equation}
     T(\tau)= T_{amb}(\tau) + 0.632( T_{ss}(\tau) - T_{amb}(\tau))
\end{equation}
en el experimento la constante de tiempo termica efectiva del sistema $\tau$ se determina al localizar el valor calculado de $T(\tau)$ en la curva de enfriamiento.
\subsection{Control proporcional integral}
Se empleó un control proporcional integral, un algoritmo de retroalimentación para regular la temperatura.
Enfriar el sistema implica una retroalimentación negativa, se define entonces el error como el calculo continuo de la diferencia entre el valor real y el valor deseado 
\begin{equation}
    e(t)= T(t)-T_{obj}(t)
\end{equation}
el termino proporcional se define como  
\begin{equation}
    p(t) = K_{p}e(t)
\end{equation}
la parte integral se encarga de sumar todos los errores que han ocurrido a lo largo del tiempo 
\begin{equation}
    Int(t) = \int_{o}^{t}e(\tau)d\tau
\end{equation}
la señal de control se define como la suma de las partes proporcional e integral, y en este caso corresponde a la corriente suministrada al sistema a lo largo  del tiempo  
\begin{equation}
    \Delta I(t) = P(t) + Int(t)
\end{equation}
\begin{equation}
   \Delta I(t) =  K_{p}e(t) + K_{I}\int_{o}^{t}e(\tau)d\tau
\end{equation}
se toma la transformada de Lapalce 
\begin{equation}
     \Lap{C(t)} =  \Lap{K_{p}e(t)} + \Lap{K_{I}\int_{o}^{t}e(\tau)d\tau}
\end{equation}
\begin{equation}
    \Delta I(s)= K_{p}e(s) + K_{I}\frac{e(s)}{s}
\end{equation}
\begin{equation}
    \Delta I(s) = e(s)(K_{p} + \frac{K_{I}}{s})
\end{equation}
se define entonces la  función de transferencia para control de corriente  
\begin{equation}
    C(s) = \frac{\Delta I(s)}{e(s)} = (K_{p} + \frac{K_{I}}{s})
\end{equation}
multiplicar (2.44) por (2.72) permite encontrar una función de transferencia para el control de la temperatura 
\begin{equation}
    C(s)G(s)= \frac{\Delta I(s)}{e(s)}\frac{\Delta T(s)}{\Delta I(s)} = \frac{\Delta T(s)}{e(s)} = (K_{p} + \frac{K_{I}}{s})\frac{-Ke^{-Ls}}{\tau s +1}
\end{equation}
se puede escribir entonces 
\begin{equation}
    C(s)G(s) = K_{p}(1+\frac{K_{I}}{K_{p}s})\frac{-Ke^{-Ls}}{\tau s +1}
\end{equation}
se consideran las unidades del cosiente $\frac{K_{I}}{K_{p}}$
\begin{equation}
    [\frac{K_{p}}{K_{I}}] = \frac{\frac{A}{C°}}{\frac{A}{C°s}} = s
\end{equation}
se define entonces este cociente en terminos de un timepo integral $t_{I}$
\begin{equation}
    t_{I} = \frac{K_{p}}{K_{I}}
\end{equation}
y se puede escribir 
\begin{equation}
    C(s)G(s)= K_{p}(1+\frac{1}{t_{I}s})\frac{-Ke^{-Ls}}{\tau s +1}
\end{equation}
\begin{equation}
 C(s)G(s)= K_{p}(\frac{t_{I}s+1}{t_{I}s})\frac{-Ke^{-Ls}}{\tau s +1}
\end{equation}

como $\tau$ tiene unidades de siempo se toma entonces
\begin{equation}
     \tau = t_{I}
\end{equation}
y (2.78) se simplifica 
\begin{equation}
    C(s)G(s) = \frac{-K_{p}Ke^{-Ls}}{\tau s}
\end{equation}
para un retardo de tiempo  muy pequeño se puede escribir la magnitud efectiva del lazo cerrado como
\begin{equation}
     C(s)G(s) = \frac{K_{p}|K|}{\tau s}
\end{equation}
\subsection{Cierre del lazo de realimentación negativa}
se conideran tres señales la de referencia r(s), el error e(s) y la salida  y(s), para una realimentacion negativa se define 
\begin{equation}
    e(s) = r(s) - y(s)
\end{equation}
de la ecuacion (2.76) se puede obtener una expresion para la salida del controlador 
\begin{equation}
    C(s)G(s)  = \frac{\Delta T(s)}{e(s)}
\end{equation}
haciendo $y(s) = \Delta T(s)$ se consigue 
\begin{equation}
    y(s)= C(s)G(s)e(s)
\end{equation}
se introduce la expresión para el error, es decir 
\begin{equation}
    y(s) = C(s)G(s)(r(s)-y(s))
\end{equation}
haciendo $C(s)G(s) = L(s)$ se puede escribir
\begin{equation}
    y(s) = L(s)r(s) - L(s)y(s)
\end{equation}
\begin{equation}
    \frac{y(s)}{r(s)} = \frac{L(s)}{1+L(s)}
\end{equation}
se obtiene entonces la respuesta de lazo cerrado H(s)
\begin{equation}
    H(s) = \frac{y(s)}{r(s)} = \frac{\frac{K_{p}|K|}{\tau s}}{1+\frac{K_{p}|K|}{\tau s}}
\end{equation}
se lleva esta expresion a una que tenga forma de un istema de primer orden 
\begin{equation}
    H(s) = \frac{K_{p}|K|}{\tau s + K_{p}|K|}
\end{equation}
\begin{equation}
    H(s) = \frac{1}{\frac{\tau}{K_{p}|K|}s + 1}
\end{equation}
  se toma 
  \begin{equation}
    \lambda = \frac{\tau}{K_{p}|K|}
  \end{equation}
   se puede escribir entonces 
   \begin{equation}
    H(s) = \frac{1}{\lambda s + 1}
   \end{equation}
   $\lambda$ es una constante que se adecua para mediar la agresividad de la respuesta de control.
   Las ecuaciones (2.79) , (2.82) y (2.94) constutuyen el puente que relaciona a las constantes del control proporcional con la caractezicacion del sistema termico.
   Este metodo se conoce como control por modelo interno, si se conoce el modelo de una planta, se puede construiir el controlador usando explicitamente ese modelo y elejir mediante la constante $\lambda$ que tan agresiva se quiere la respuesta. 
   